% !TeX root = ../../Book1.tex
% 纲要标题：## 附录 E　有序域与实闭域
% 纲要位置：计划纲要.md，第 910 行。
% BEGIN APPENDIX OUTLINE SYNC
% > 域论后的短篇接口，不加入基础 Galois 理论的先修链。这里的“实代数”指序、实闭性和正性相关的代数理论，不只是“以实数为标量的代数”。资料依据见修订说明 [R1]–[R2]。
% END APPENDIX OUTLINE SYNC

\chapter{有序域与实闭域}
\label{b1:app:E}

\begin{chapterabstract}
有序域在四则运算之外增加正负比较。形式实性用平方和刻画一个域能否带序；实闭性则给出实数域在多项式问题上的代数特征。本附录证明可赋序判据、多项式中间值性质和 Sturm 根计数，借助明确列出的实闭域结构定理说明这些结论为何不依赖实数的上确界性质。
\end{chapterabstract}

先修是域扩张、多项式除法与形式导数。可赋序判据使用\cref{b1:thm:A-zorn}。实闭域的等价刻画与给定序的实闭包存在唯一性作为外部结构结果采用，并在陈述旁标明来源；本附录不声称补齐它们的全部证明。这里的实代数研究序、实闭性与正性，不只是把标量域选为实数。它不加入基础 Galois 理论的先修链。

% 纲要内容（仅作写作计划，尚非教材正文）：
%
% > 域论后的短篇接口，不加入基础 Galois 理论的先修链。这里的“实代数”指序、实闭性和正性相关的代数理论，不只是“以实数为标量的代数”。资料依据见修订说明 [R1]–[R2]。
%

% !TeX root = ../../Book1.tex
% 纲要标题：### E.1 域上的序与形式实性
% 纲要位置：计划纲要.md，第 914 行。
% BEGIN APPENDIX OUTLINE SYNC
% #### E.1.1 有序域、正锥与平方和
%
% #### E.1.2 形式实域与 Artin–Schreier 判据
%
% #### E.1.3 域上的序与二次型
% END APPENDIX OUTLINE SYNC

\section{域上的序与形式实性}
\label{b1:sec:E01}

说一个数为正，需要与加法、乘法相容的比较规则。域公理本身不含这样的规则：复数域不能带域序，而同一个有理函数域却可以有许多不同的序。本节先给出正锥的代数条件，再将可赋序性转成平方和问题。

% 纲要内容（仅作写作计划，尚非教材正文）：
%
% * 有序域、正锥与平方和；实数域和有理函数域的例子
% * 形式实域的定义：\(-1\) 不是有限个平方之和；Artin–Schreier 可赋序判据
% * 指定一个序与枚举域上的全部序是不同问题；与第二卷二次型的联系
%

\subsection{有序域、正锥与平方和}
\label{b1:subsec:E01:01}

\begin{definition}\label{b1:def:ordered-field}
\term[youxu-yu]{有序域}[ordered field]是域 \(K\) 连同全序 \(\leq\)，使 \(a\leq b\) 蕴含 \(a+c\leq b+c\)，并使 \(a,b\geq0\) 蕴含 \(ab\geq0\)。它的\term[zheng-zhui]{正锥}[positive cone]在本附录采用含零约定，即 \(P=\{a\in K\mid a\geq0\}\)。
\end{definition}

正锥的等价描述为
\[
P+P\subseteq P,\qquad PP\subseteq P,\qquad
P\cup(-P)=K,\qquad P\cap(-P)=\{0\}.
\]
由这些条件反过来定义 \(a\leq b\) 当且仅当 \(b-a\in P\)，可分别检验自反性、传递性、反对称性、全序性与运算相容性。非零平方总为正：非零 \(a\) 与 \(-a\) 必有一个为正，而两者的平方相同。特别 \(1>0\)，所以所有正整数的像为正；有序域必有特征零。

\ProseParagraphBreak
通常的 \(\mathbb Q\) 和 \(\mathbb R\) 是有序域。对 \(\mathbb Q(t)\)，规定非零有理函数 \(f/g\) 为正当且仅当 \(f,g\) 的最高次项系数同号。这与表示无关，因为换表示时交叉乘积相等；乘法保持正号，加法可先通分到正的平方分母，两个正最高次项不会相消。因此它给出一个序，且 \(t>n\) 对所有整数 \(n\) 成立。以 \(t\mapsto-t\) 拉回这个序，则得到 \(-t>n\) 的另一个序。

\begin{definition}\label{b1:def:sum-of-squares}
域 \(K\) 的平方和集合记为
\[
\Sigma K^2=\left\{\sum_{j=1}^r a_j^2\ \middle|\ r\geq0,\ a_j\in K\right\};
\]
空和为零。若 \(-1\notin\Sigma K^2\)，称 \(K\) 为\term[xingshi-shiyu]{形式实域}[formally real field]。
\end{definition}
\symidx{\(\Sigma K^2\)：域中有限个平方的和}

任一域序都使平方和非负，所以可赋序的域必形式实。形式实性只涉及四则运算，没有预先选择哪个元素为正。有限域不满足它，因为特征 \(p\) 时 \(-1\) 是 \(p-1\) 个 \(1^2\) 的和；\(\mathbb C\) 中则已有 \(-1=i^2\)。

\subsection{形式实域与 Artin–Schreier 判据}
\label{b1:subsec:E01:02}

下面的\term[Artin-Schreier-kafuxu-panju]{Artin--Schreier 可赋序判据}[Artin--Schreier orderability criterion]把上一段的必要条件变成充分条件。它与\secref{b1:sec:2004}中正特征循环扩张的 Artin--Schreier 理论属于不同的问题。

\begin{theorem}[Artin--Schreier]\label{b1:thm:artin-schreier-ordering}
域 \(K\) 能够赋予一个域序，当且仅当它形式实。
\end{theorem}
\begin{proof}
必要性已经说明。设 \(K\) 形式实。考虑所有子集 \(T\subseteq K\)，要求它含有全部平方，对加法、乘法封闭，且不含 \(-1\)。这一集合族非空，因为 \(\Sigma K^2\) 满足条件。按包含关系取任一链，其并仍含全部平方且不含 \(-1\)；加、乘涉及有限个元素，能在链中同一个成员里检验，所以并仍满足条件。由\cref{b1:thm:A-zorn}，存在极大成员 \(P\)。

若 \(s\in P\) 且 \(s\neq0\)，则 \(s^{-1}=s(s^{-1})^2\in P\)。所以 \(a,-a\in P\) 且 \(a\neq0\) 将推出 \(-1=(-a)a^{-1}\in P\)，不可能。因此 \(P\cap(-P)=\{0\}\)。

还需证明每个 \(a\) 或其负元属于 \(P\)。若 \(a\notin P\)，考察
\[
T=P-aP=\{r-as\mid r,s\in P\}.
\]
它包含 \(P\)，对加法封闭；乘法展开中的 \(a^2\) 属于 \(P\)，故也对乘法封闭。若 \(-1=r-as\)，则 \(s\neq0\)，从而 \(a=(1+r)s^{-1}\in P\)，矛盾。所以 \(-1\notin T\)，极大性迫使 \(T=P\)。取 \(r=0,s=1\) 得 \(-a\in P\)。至此正锥的全部条件成立，\cref{b1:def:ordered-field}后的构造给出所需域序。
\end{proof}

这个证明还说明：若一个集合已经含全部平方、对加乘封闭且不含 \(-1\)，就能扩成正锥。这里的极大性只用来完成符号选择，不是计算某个具体序的有限算法。

\subsection{域上的序与二次型}
\label{b1:subsec:E01:03}

指定一个序以后，所有非零平方都为正，但非平方元的符号可能随序改变。例如 \(\mathbb Q(\sqrt2)\) 到 \(\mathbb R\) 的两个嵌入把 \(\sqrt2\) 分别送到正根、负根，拉回两个不同的序。可赋序判据只保证至少存在一个序，不枚举所有序。

\ProseParagraphBreak
平方和也把这一问题连到二次型。对角二次型 \(q(x_1,\ldots,x_n)=a_1x_1^2+\cdots+a_nx_n^2\) 的系数在某个固定序中全为正时，任何非零向量都使 \(q\) 为正。例如 \(x^2+y^2=0\) 在有序域中只能有零解；若允许 \(-1\) 为平方，结论立即失败。更一般地，形式实性等价于任意平方和为零时各项均为零，因为只要有一个非零项，就可移项并除以它的平方得到 \(-1\) 的平方和表示。第二卷对二次型作系统研究时，还要区分等价、各向同性与所选域序；本节没有给出一般二次型的分类。

% !TeX root = ../../Book1.tex
% 纲要标题：### E.2 实闭域与实闭包
% 纲要位置：计划纲要.md，第 922 行。
% BEGIN APPENDIX OUTLINE SYNC
% #### E.2.1 实闭域的等价刻画
% #### E.2.2 实闭包的存在性与唯一性
% #### E.2.3 实闭性与序完备性
% END APPENDIX OUTLINE SYNC

\section{实闭域与实闭包}
\label{b1:sec:E02}

实数域并不代数闭：\(X^2+1\) 无实根。但在代数扩张中再加入任何缺少的根，便不能一直保留实数意义的正负比较。实闭性抽取这种代数上的极大性质，所需的结构定理与实数的序完备性必须分清。

% 纲要内容（仅作写作计划，尚非教材正文）：
%
% #### E.2.1 实闭域的等价刻画
% #### E.2.2 实闭包的存在性与唯一性
% #### E.2.3 实闭性与序完备性
%

\subsection{实闭域的等价刻画}
\label{b1:subsec:E02:01}

\begin{definition}\label{b1:def:real-closed-field}
若域 \(R\) 形式实，且没有真代数扩张仍为形式实域，则称 \(R\) 为\term[shibi-yu]{实闭域}[real closed field]。
\end{definition}

这个定义同时要求两件事：域本身可以赋序，而在它上面再作任何真代数扩张，都不可能赋予域序。限定“代数”很关键，因为这里讨论的是加入多项式的根，不是任意增加新的元素。例如 \(\mathbb Q\) 形式实，但 \(\mathbb Q(\sqrt2)\) 仍可以沿实数的序比较大小，所以 \(\mathbb Q\) 不实闭。“形式实”仅排除 \(-1\) 为平方和，“实闭”还要求这种性质不能在更大的代数域中继续保持。

\ProseParagraphBreak
以下定理把实闭域的极大性质改写成关于多项式的条件；后面的根计数证明将用到这些条件。等价刻画亦见\cite[Theorem A.5.19]{Mangolte2020RealAlgebraicVarieties}。
\begin{theorem}[实闭域的等价刻画]\label{b1:thm:real-closed-characterizations}
对形式实域 \(R\)，下列条件等价：
\begin{enumerate}
\item \(R\) 实闭；
\item \(R\) 有一个域序，使每个正元素都是平方，且 \(R[X]\) 中每个奇数次多项式都在 \(R\) 中有根；
\item 取 \(i^2=-1\)，二次扩张 \(R(i)\) 是代数闭域。
\end{enumerate}

这些条件成立时，\(R\) 的域序唯一，其正元素恰为非零平方。
\end{theorem}

\begin{proof}
先证第一项推出第二项。由\cref{b1:thm:artin-schreier-ordering}，给 \(R\) 取一个域序。若正元素 \(a\) 不是平方，则在二次扩张 \(R(\sqrt a)\) 中，等式 \(-1=\sum_j(x_j+y_j\sqrt a)^2\) 的常数项会给出
\[
-1=\sum_jx_j^2+a\sum_jy_j^2\ge0,
\]
矛盾。因此该扩张形式实，与实闭性矛盾，故每个正元素都是平方。

若有奇数次多项式无根，取次数大于一且最小的奇数次不可约多项式 \(f\)，记其次数为 \(n\)。域 \(R[X]/(f)\) 不是形式实，所以有次数均小于 \(n\) 的 \(g_j\)，使
\[
1+\sum_jg_j(X)^2=f(X)h(X).
\]
左侧的最高次系数为非零平方和，不会抵消，故其次数为偶数且至多 \(2n-2\)。于是 \(h\) 为奇数次多项式，次数小于 \(n\)。由 \(n\) 的最小性，\(h\) 的某个奇数次不可约因子必为一次，故 \(h\) 有根 \(b\in R\)。代入上式又得 \(1+\sum_jg_j(b)^2=0\)，矛盾。

再证第二项推出第三项。令 \(F=R(i)\)。\(F\) 中每个元素都是平方：对 \(a+bi\)，若 \(b\ne0\)，取正平方根 \(r=\sqrt{a^2+b^2}\)、\(u=\sqrt{(r+a)/2}\)，则 \(r>|a|\)、\(u>0\)，并有
\[
\left(u+\frac b{2u}i\right)^2=a+bi.
\]
若 \(b=0\)，按 \(a\) 的正负分别用实平方根或 \(i\sqrt{-a}\)，零元另取零。因而 \(F\) 没有二次扩张。

任取 \(F\) 上的一个多项式，令 \(N/R\) 为包含它的全部根及 \(F\) 的有限正规扩张。特征为零，故 \(N/R\) 为 Galois 扩张。取 \(G=\operatorname{Gal}(N/R)\) 的一个 Sylow 二子群 \(P\)。不动域 \(N^P/R\) 的次数为奇数；由\cref{b1:thm:primitive-element}，若这个扩张非平凡，就有次数大于一的奇数次不可约多项式，与第二项矛盾。因此 \(N^P=R\)、\(G=P\) 是二群。非平凡有限二群 \(H\) 总有指数二的子群：取极大真子群 \(M\)，\cref{b1:prop:p-normalizer-condition}使 \(M\trianglelefteq H\)，商若阶数大于二，Cauchy 定理给出它的非平凡真子群，与 \(M\) 极大矛盾。若 \(\operatorname{Gal}(N/F)\ne1\)，将此事实和 Galois 对应合用，便给出 \(F\) 的二次扩张，仍矛盾。故 \(N=F\)，原多项式的全部根已经在 \(F\) 中。这证明 \(F\) 代数闭。

第三项推出第一项时，若 \(L/R\) 为代数扩张，则\cref{b1:thm:embedding-extension}使它可嵌入代数闭域 \(R(i)\)。这个二次扩张没有其他中间域，所以 \(L\) 若为真扩张，便同构于 \(R(i)\)，包含 \(-1\) 的平方根而非形式实。因此 \(R\) 实闭。

最后，第二项说明正元素恰为非零平方，任何域序都必须把非零平方判为正，因此序唯一。
\end{proof}

三种条件分别适合不同的问题。第一项从扩张出发，说明不能再加入代数元而保持可赋序性；第二项允许直接检验两类方程，即正元素的平方根和奇数次多项式的根；第三项则说明只需加入一个 \(-1\) 的平方根，就足以使所有多项式分裂。第二项并未要求所有偶数次多项式有根，例如 \(X^2+1\) 始终无根，也未要求奇数次多项式只有一个根，例如 \(X^3-X\) 有三个根。

\ProseParagraphBreak
序唯一性意味着所有正负都由域运算决定。条件三也说明 \(R\) 自身并不代数闭，因为形式实性排除了 \(i\in R\)。

\begin{corollary}\label{b1:cor:real-polynomial-factorization}
实闭域 \(R\) 上的每个非零多项式都是非零常数、一次因子和形如 \((X-a)^2+b^2\)、\(a,b\in R\)、\(b\neq0\) 的二次因子的乘积。后面这些首一二次因子在 \(R\) 上处处为正。
\end{corollary}

\begin{proof}
由\cref{b1:thm:real-closed-characterizations}，多项式在 \(R(i)\) 中完全分裂。共轭映射 \(a+bi\mapsto a-bi\) 固定系数，所以非实根按共轭对出现，重数也相同。每对对应
\[
\bigl(X-(a+bi)\bigr)\bigl(X-(a-bi)\bigr)=(X-a)^2+b^2.
\]
实根给出一次因子。平方非负而 \(b^2>0\)，所以每个二次因子处处为正。
\end{proof}

因此实闭域上的不可约多项式只有一次和二次两类。没有实根的多项式未必不可约，例如 \(X^4+3X^2+2=(X^2+1)(X^2+2)\) 无实根，却分成两个二次因子。后面判断符号时真正使用的是：每个非实根及其共轭合成的首一二次因子永远为正，故只有一次因子能够让符号改变。

\subsection{实闭包的存在性与唯一性}
\label{b1:subsec:E02:02}

\begin{definition}\label{b1:def:real-closure}
设 \((K,\leq)\) 为有序域。若 \(R\) 是实闭域，\(K\hookrightarrow R\) 是保持序的域嵌入，且 \(R/K\) 为代数扩张，则称 \(R\) 连同该嵌入为 \((K,\leq)\) 的一个\term[shibi-bao]{实闭包}[real closure]。
\end{definition}

保持序的意思是：若 \(a,b\in K\) 且 \(a<b\)，它们在 \(R\) 中仍满足同一不等式。代数性则要求 \(R\) 的每个元素都满足一个系数在 \(K\) 中的非零多项式。寻找实闭包，就是在保留这些原有符号的前提下补足代数根，直到得到实闭域。不能直接使用代数闭包来代替，因为代数闭包含有 \(-1\) 的平方根，已经不能赋序。

\ProseParagraphBreak
下述定理固定的是有序域，不仅是遗忘了序的域。原始论证见\cite[Satz 8]{ArtinSchreier1927RealFields}。
\begin{theorem}[实闭包]\label{b1:thm:real-closure-existence}
每个有序域都有实闭包。若 \(R_1,R_2\) 是同一个有序域 \(K\) 的两个实闭包，则存在唯一保持序的 \(K\)-同构 \(R_1\rightarrow R_2\)。
\end{theorem}

\begin{proofsketch}
存在性可直接由极大有序代数扩张构造。固定 \(K\) 的代数闭包 \(\Omega\)，考虑其中包含 \(K\) 的子域及其延伸给定序的正锥，按子域包含和正锥延伸排序。链的子域与正锥分别取并，仍为有序代数扩张；Zorn 引理给出极大者 \(R\)。

核对极大者实闭，需要两项保持序的扩张事实。正元素 \(a\) 的平方根可以加入而保留序：否则由\cref{b1:thm:artin-schreier-ordering}证明末的正锥扩充方法，已有正锥生成的预序中会出现
\[
-1=\sum_jp_j(x_j+y_j\sqrt a)^2,\qquad p_j\ge0.
\]
比较常数项得到 \(-1=\sum_jp_j(x_j^2+ay_j^2)\ge0\)，矛盾。奇数次扩张同样保留序。若不然，取次数 \(n\) 最小的反例，用本原元写成 \(E[X]/(f)\)，其中 \(E\) 是所考虑的有序基域，\(n=\deg f\) 为奇数。不能扩充正锥意味着存在 \(p_j\ge0\)、\(\deg g_j<n\)，使
\[
1+\sum_jp_jg_j(X)^2=f(X)h(X).
\]
最高次项为正，故 \(\deg h\) 为小于 \(n\) 的奇数。取 \(h\) 的一个奇数次不可约因子及其根 \(\beta\)，代入上式，就在更小的奇数次扩张中得到同样的 \(-1\) 表示，与最小性矛盾。因此这两类扩张都能保留指定序。极大性使 \(R\) 的正元素已有平方根，奇数次多项式已有根；上一条等价刻画说明 \(R\) 实闭。

同构的存在性使用一个代数根的符号判定事实：给定有序域 \(E\)、不可约 \(f\in E[X]\) 以及它在实闭扩张中的第 \(j\) 个实根 \(\alpha\)，每个 \(g(\alpha)\)、\(g\in E[X]\)，的符号由 \(E\) 的序及这个根的序号决定。其证明对 \(f,g\) 和导数作带符号的 Euclid 余式计算，先数实根，再判定各根上的符号；遇到余式为零时降低次数继续，故只用有限次域运算和系数符号比较。这也是后一节 Sturm 方法的根符号判定推广。这里省略这项推广的余式表与各退化情形，实闭包的原始唯一性论证见\cite[Satz 8]{ArtinSchreier1927RealFields}，根符号判定和消去法见\cite[Sections 1.2, 1.4]{BochnakCosteRoy1998RealAG}。

现在取所有从 \(R_1\) 的中间域到 \(R_2\) 的保持序的部分 \(K\)-嵌入。链的并给出嵌入，取极大者 \(\phi:E\to R_2\)。若 \(\alpha\in R_1\setminus E\)，它的最小多项式经 \(\phi\) 变换后，在 \(R_2\) 中有同样多的实根；将 \(\alpha\) 送到相同序号的根。根符号判定使 \(g(\alpha)\) 的符号全部保持，因而这个域嵌入可保持序地延伸到 \(E(\alpha)\)，与极大性矛盾。故定义域为全部 \(R_1\)。像为实闭域，\(R_2\) 又在此像上代数且形式实，实闭性的极大性质使像等于 \(R_2\)。这样得到同构。

最后，保持序的 \(K\)-自同构会置换任意 \(f\in K[X]\) 的有限个实根，同时保持它们的次序，因此逐个固定。\(R_1/K\) 的代数性使它固定所有元素，故只能是恒等。把一个同构的逆与另一个复合，得到这样的自同构，因而同构唯一。
\end{proofsketch}

代数性不能删去。例如 \(\mathbb R\) 是包含 \(\mathbb Q\) 的实闭域，但不是 \(\mathbb Q\) 的实闭包。另一个容易混淆的条件是保持给定序：对 \(\mathbb Q(t)\)，取 \(t>0\) 的实闭包会使 \(t\) 成为平方，取 \(t<0\) 的实闭包则不可能如此，所以两者不能通过固定 \(t\) 的域同构相连。

\ProseParagraphBreak
唯一性中的“唯一同构”比一般代数闭包的唯一性更强，原因是保持序的同构不能互换一组有序的实根。

\subsection{实闭性与序完备性}
\label{b1:subsec:E02:03}

设 \((K,\leq)\) 为有序域。若 \(K\) 的每个非空且有上界的子集都在 \(K\) 中有上确界，即有一个最小上界，则称 \((K,\leq)\) 具有\term[xu-wanbeixing]{序完备性}[order completeness]。这个条件要求处理任意子集；实闭性则只规定有限个系数组成的多项式及其代数根。两个性质的要求不同，下面用具体域比较它们。

\ProseParagraphBreak
\(\mathbb R\) 的正元素有平方根，奇数次实多项式由通常的中间值定理有根，所以\cref{b1:thm:real-closed-characterizations}说明它实闭。此处只是利用实分析已经建立的性质来识别这个例子；下面从实闭性证明一般多项式的中间值性质，不以实数完备性作为一般有序域的公理。

\begin{proposition}\label{b1:prop:real-algebraic-closure}
实代数数域 \(\mathbb R_{\mathrm{alg}}=\{x\in\mathbb R\mid x\text{ 在 }\mathbb Q\text{ 上代数}\}\) 是 \(\mathbb Q\) 在通常序下的实闭包。
\end{proposition}

\begin{proof}
代数元在加、减、乘、除下封闭，所以这是域。其正元素的正平方根仍在 \(\mathbb R\) 中，并对原元素为二次代数，因此在 \(\mathbb Q\) 上代数。对一个系数在此域内的奇数次多项式，先在 \(\mathbb R\) 中取实根。有限个系数生成 \(\mathbb Q\) 的有限扩张，根对该有限扩张代数，由塔式公式在 \(\mathbb Q\) 上也代数。因此它满足\cref{b1:thm:real-closed-characterizations}的第二项，故实闭；它按定义在 \(\mathbb Q\) 上代数，且序与 \(\mathbb Q\) 相容，因而为实闭包。
\end{proof}

\(\mathbb R_{\mathrm{alg}}\) 没有实数域的序完备性。取任一实超越数 \(\tau\)，考虑子集 \(S=\{a\in\mathbb R_{\mathrm{alg}}\mid a<\tau\}\)。取比 \(\tau\) 小及大的整数，分别说明 \(S\) 非空并且有上界。若某个 \(s\in\mathbb R_{\mathrm{alg}}\) 是它的上确界，则 \(s\neq\tau\)。当 \(s<\tau\) 时，由有理数在实数中的稠密性，存在 \(q\in\mathbb Q\) 满足 \(s<q<\tau\)，于是 \(q\in S\)，说明 \(s\) 甚至不是上界。当 \(s>\tau\) 时，取 \(\tau<q<s\) 的有理数，便得到比 \(s\) 更小的上界。因此任何 \(s\) 都不可能是上确界。

\ProseParagraphBreak
设 \((K,\leq)\) 为有序域。若对每个 \(x\in K\)，都存在正整数 \(n\) 使 \(x<n\)，则称 \((K,\leq)\) 为\term[Archimedean-youxu-yu]{Archimedean 有序域}[Archimedean ordered field]。\(\mathbb Q\)、\(\mathbb R_{\mathrm{alg}}\) 与 \(\mathbb R\) 都满足这一性质。还可以取\secref{b1:sec:E01}中使 \(t>n\) 对所有正整数成立的 \(\mathbb Q(t)\) 的实闭包。由\cref{b1:thm:real-closure-existence}，这个实闭包存在且保留给定序，所以原不等式全部保留；其中的 \(t\) 仍大于每个正整数。因此实闭性既不蕴含 Archimedean 性，也不蕴含序完备性。

% !TeX root = ../../Book1.tex
% 纲要标题：### E.3 多项式的实根与实代数
% 纲要位置：计划纲要.md，第 930 行。
% BEGIN APPENDIX OUTLINE SYNC
% #### E.3.1 中间值性质与 Sturm 根计数
% #### E.3.2 传递原理、量词消去与半代数集合
% ---
% END APPENDIX OUTLINE SYNC

\section{多项式的实根与实代数}
\label{b1:sec:E03}

实根计数可以只用多项式除法与符号比较完成。本节先从实闭域上的因式分解推出中间值性质，再给出 Sturm 序列的计数公式。即使域中有大于所有整数的元素，证明也仍然适用。

% 纲要内容（仅作写作计划，尚非教材正文）：
%
% #### E.3.1 中间值性质与 Sturm 根计数
% #### E.3.2 传递原理、量词消去与半代数集合
% ---
%

\subsection{中间值性质与 Sturm 根计数}
\label{b1:subsec:E03:01}

\begin{theorem}[多项式中间值性质]\label{b1:thm:real-closed-ivt}
设 \(R\) 实闭，\(a<b\) 属于 \(R\)，且 \(f\in R[X]\) 满足 \(f(a)f(b)<0\)。则存在 \(c\in R\)，使 \(a<c<b\) 且 \(f(c)=0\)。
\end{theorem}

\begin{proof}
由\cref{b1:cor:real-polynomial-factorization}，将 \(f\) 分成常数、一次因子和处处为正的二次因子。若没有实根位于 \((a,b)\)，每个一次因子在两端的符号相同，二次因子在两端均为正，故乘积 \(f(a)f(b)>0\)，矛盾。
\end{proof}

例如，在实闭域中 \(f=X^3-X-1\) 满足 \(f(1)=-1\)、\(f(2)=5\)，所以 \((1,2)\) 中至少有一个根。仅看端点异号还不知道有几个根；如果两端同号，也不能断言无根，例如 \((X-1)(X-2)\) 在 \(0\) 和 \(3\) 处都为正，却在中间有两个根。

\ProseParagraphBreak
因式分解还说明：非零多项式的符号在不含根的区间上固定。根计数因此可以转化为研究越过各个根时符号怎样改变。Sturm 方法把 \(f\)、\(f'\) 和一串负余数放在一起：\(f'\) 帮助识别 \(f\) 的简单根，而余数前的负号使其他多项式的零点不会改变总变号数。以下把这两个作用分别证明。

\ProseParagraphBreak
先设 \(f\in R[X]\) 为无重根的非常数多项式。它与 \(f'\) 互素：若某个因子同时整除两者，就会给出重根。令 \(f_0=f\)、\(f_1=f'\)，将 \(f_{j-1}\) 除以 \(f_j\)，把所得非零余数的相反数记为 \(f_{j+1}\)，即
\begin{equation}\label{b1:eq:sturm-remainders}
f_{j-1}=q_jf_j-f_{j+1},\qquad \deg f_{j+1}<\deg f_j,
\end{equation}
如果余数为零便停止，记最后一个非零多项式为 \(f_m\)。余数次数严格下降保证过程终止；Euclid 算法说明最后一项与 \(\gcd(f,f')\) 相差一个非零常数因子，所以互素性保证 \(f_m\) 是非零常数。所得序列称为 \(f\) 的\term[Sturm-xulie]{Sturm 序列}[Sturm sequence]。对 \(x\in R\)，把 \(f_0(x),\ldots,f_m(x)\) 中的零项删去，记剩余符号列的变号次数为 \(V(x)\)：每当相邻两个非零值的符号相反，就计一次。例如符号列 \((+,0,-,-,+)\) 先删成 \((+,-,-,+)\)，共有两次变号，而不是把零另当作第三种符号。
\symidx{\(V(x)\)：Sturm 序列在一点的变号次数}

\begin{theorem}[Sturm]\label{b1:thm:sturm-count}
在上述条件下，若 \(a<b\) 且 \(f(a)f(b)\neq0\)，则 \(f\) 在 \((a,b)\) 中的根数为 \(V(a)-V(b)\)，每个根只计一次。
\end{theorem}

\begin{proof}
相邻两个 \(f_j\) 不可能有公共根。若 \(f_j(c)=f_{j+1}(c)=0\)，把递推式逐项向后使用，会迫使 \(f_{j+2}(c),\ldots,f_m(c)\) 都为零，最终与 \(f_m\) 是非零常数矛盾。

先考察一个中间项的根，即 \(f_j(c)=0\)、\(1\leq j<m\)。递推式给出 \(f_{j-1}(c)=-f_{j+1}(c)\neq0\)。取 \(c\) 两侧足够近的位置，使两个邻项都不经过自己的根，则其符号仍相反。假如邻项分别是正、负，中间项可能使三项成为 \((+,+,-)\) 或 \((+,-,-)\)，两者都只贡献一次变号；邻项为负、正时结论相同。在 \(c\) 处删去零项后，两邻项也恰贡献一次。因此经过中间项的零点不改变总变号数。若多个中间项在 \(c\) 处为零，它们不相邻；各自涉及的相邻符号对不重复，所以上述计算可以同时使用。

再考察 \(f\) 的根 \(c\)。无重根性给出 \(f'(c)\neq0\)。写成 \(f(X)=(X-c)g(X)\)，求导并代入 \(c\) 得 \(g(c)=f'(c)\neq0\)。在不跨越 \(g\) 和 \(f'\) 其他根的小区间里，两者符号不变且相同。\(c\) 左侧的 \(X-c\) 为负，所以 \(f\) 与 \(f'\) 异号；右侧的 \(X-c\) 为正，所以两者同号。因此从左到右经过 \(c\)，最前一对的变号次数由一变成零，恰减少一；后面各项按前段不改变总变号数。

这些局部比较并不需要分析学中的连续性定理。每个非零多项式只有有限个根，将全部 \(f_j\) 的根按序排列，相邻根之间的各项符号由\cref{b1:cor:real-polynomial-factorization}保持不变。在需要选择比较点时，取两个相邻根的中点即可；在端点外侧也可以取半个正距离。由此从 \(a\) 到 \(b\) 逐段比较，变号数只在 \(f\) 的根处各减少一。端点若有中间项为零，前段的删零规则保证 \(V(a)\)、\(V(b)\) 与邻近区间上的计数相同。总减少量因而恰为 \((a,b)\) 内 \(f\) 的根数。
\end{proof}

若原多项式有重根，先以 \(h=f/\gcd(f,f')\) 代替 \(f\)。实闭域特征为零；对每个不可约因子比较导数中的重数，可知 \(h\) 恰保留各因子一次，因此与 \(f\) 有相同的互异根，且无重根。上述算法于是计算互异实根数；若要计算重数，另保留原分解中的指数。

\ProseParagraphBreak
回到 \(f=X^3-X-1\)。先取 \(f_1=3X^2-1\)，逐次相除得
\[
\begin{aligned}
X^3-X-1&=\frac X3(3X^2-1)-\left(\frac23X+1\right),\\
3X^2-1&=\left(\frac92X-\frac{27}4\right)
\left(\frac23X+1\right)+\frac{23}4.
\end{aligned}
\]
第二式的余数是 \(23/4\)，其相反数才是最后的 Sturm 项。因此得到
\[
f_0=X^3-X-1,\quad f_1=3X^2-1,\quad
f_2=\tfrac23X+1,\quad f_3=-\tfrac{23}{4}.
\]
在 \(1\) 处四个值为 \(-1,2,5/3,-23/4\)，符号列是 \((-,+,+,-)\)，所以 \(V(1)=2\)。在 \(2\) 处四个值为 \(5,11,7/3,-23/4\)，符号列是 \((+,+,+,-)\)，所以 \(V(2)=1\)。于是 \((1,2)\) 内恰有一个实根，这比中间值性质提供的“至少一个”更精确。

\ProseParagraphBreak
同样，在 \(-2\) 处符号列为 \((-,+,-,-)\)，所以 \(V(-2)=2\)，\((-2,2)\) 中也只包含这一个根。当 \(x\geq2\) 时，\(x(x^2-1)-1\geq5\)；当 \(x\leq-2\) 时，该式不大于 \(-7\)。故外侧没有根，已经数出了全部实根。以上仅用有理数运算和符号比较，在每个实闭域中都成立。一般使用 Sturm 方法时，若工作域中的系数能精确运算并能判定符号，便可依此执行算法；抽象域的定义本身不保证存在可用的机器表示。

\subsection{传递原理、量词消去与半代数集合}
\label{b1:subsec:E03:02}

一个\term[bandaishu-jihe]{半代数集合}[semialgebraic set]是实闭域 \(R\) 的 \(R^n\) 中，由有限多个多项式等式与严格不等式通过有限次并、交、补得到的集合。系数所在的域须随问题指定；例如圆盘可写为 \(x^2+y^2\leq1\)，其中非严格不等式可拆成严格不等式与等式。

\ProseParagraphBreak
所谓投影，是只保留一些坐标而忘去其余坐标。若 \(S\subseteq R^{n+1}\) 由条件 \(\Phi(x_1,\ldots,x_n,y)\) 描述，那么其前 \(n\) 个坐标的投影由“存在 \(y\) 使 \(\Phi\) 成立”描述。消去这个存在量词，就是把条件改写成只涉及 \(x_1,\ldots,x_n\) 的有限多个等式与不等式。例如
\[
\exists y\ (y^2=x)\quad\Longleftrightarrow\quad x\geq0
\]
在实闭域中成立：平方非负给出必要性，正元素有平方根给出充分性。对一般有限多项式条件也能消去量词，是实闭域的\term[liangci-xiaoqu]{量词消去}[quantifier elimination]定理；这里仅预告其内容，不证明一般消去算法，也不把它当作上述 Sturm 证明的前提。

\ProseParagraphBreak
这里的一阶公式由域元素之间的等式、不等式，经过“且、或、非”以及对域元素的“存在”“任意”量词组成。被量词约束的是域元素，而不是任意子集；未被量词约束的变量称为自由变量。没有自由变量的公式是句子，例如“每个正元素都有平方根”可以写成
\[
\forall x\,\bigl(x>0\ \Longrightarrow\ \exists y\ (y^2=x)\bigr).
\]
若公式还点名使用某个给定域元素，例如 \(x^2>a\) 中事先指定的 \(a\)，就称使用了参数 \(a\)。这些说明只为读懂下一项预告，不要求先学习完整的逻辑理论。

\ProseParagraphBreak
相应的\term[chuan-di-yuanli]{传递原理}[transfer principle]说：用 \(0,1,+,-,\cdot,<\) 表述、没有额外参数的一阶句子，在所有实闭域中真值相同；若实闭域 \(R\subseteq S\)，则带 \(R\) 中参数的一阶公式，在取自 \(R\) 的自由变量值上也有相同真值。这同样作为进一步理论的预告，不在此证明。它不表示任意关于实数的分析陈述都能转移。例如“每个非空有上界的子集都有上确界”量化的是任意子集，不是这里的一阶域语言，而\secref{b1:sec:E02}已给出该性质失效的实闭域。

\ProseParagraphBreak
半代数集合的投影、量词消去和正性证书将转第三卷附录 G；相应的 Tarski--Seidenberg 原理、投影定理与正性定理可参阅 Bochnak、Coste 与 Roy 的《Real Algebraic Geometry》\cite[Sections 1.4, 2.2, 4.4]{BochnakCosteRoy1998RealAG}。这里不展开模型论、赋值理论或二次型分类。读者若只需要核验一元多项式的实根数，前一小节的除法表、端点符号和\cref{b1:thm:sturm-count}已经组成独立可核验的计算过程。


\begin{chaptersummary}
形式实性等价于能够选择一个与域运算相容的序，但这个序一般不唯一。实闭域的正元素恰为非零平方，因此序由域运算决定；固定有序基域后，实闭包也有相应的唯一性。实闭域上多项式的因式分解足以支持中间值性质与 Sturm 计数，不需要假定所有有界集合都有上确界。
\end{chaptersummary}

\BookExercises

\begin{exercise}\label{b1:ex:E-two-orders}
证明 \(\mathbb Q(\sqrt2)\) 恰有两个域序，并分别说明 \(\sqrt2\) 的符号。
\end{exercise}
\begin{solution}
域到 \(\mathbb R\) 的两个嵌入 \(\sqrt2\mapsto\sqrt2\) 与 \(\sqrt2\mapsto-\sqrt2\) 拉回两个不同的序。在任一域序中，\(\mathbb Q\) 具有通常的序，\(\sqrt2\) 的符号只能正或负。若它为正，对正有理数 \(r\)，比较 \(\sqrt2\) 与 \(r\) 等价于比较 \(2\) 与 \(r^2\)，因为正元素平方严格递增。这决定所有 \(a+b\sqrt2\) 的符号：\(b=0\) 时看 \(a\)，\(b\neq0\) 时化为与有理数 \(-a/b\) 比较。负号情况用 \(-\sqrt2\) 同样处理，所以没有第三种序。
\end{solution}

\begin{exercise}\label{b1:ex:E-sum-squares-zero}
证明域 \(K\) 形式实，当且仅当每个等式 \(x_1^2+\cdots+x_n^2=0\) 都迫使所有 \(x_i=0\)。据此说明有限域和代数闭域不能赋序。
\end{exercise}
\begin{solution}
若某个 \(x_j\neq0\)，将等式移项并除以 \(x_j^2\)，得到 \(-1=\sum_{i\neq j}(x_i/x_j)^2\)，与形式实性矛盾。反向若 \(-1=\sum y_i^2\)，则 \(1^2+\sum y_i^2=0\) 给出非零解。特征 \(p\) 的域中 \(-1\) 是 \(p-1\) 个 \(1^2\) 之和；代数闭域中 \(-1\) 本身有平方根。故两类域均非形式实，由\cref{b1:thm:artin-schreier-ordering}不能赋序。
\end{solution}

\begin{exercise}\label{b1:ex:E-quadratic-positive}
设 \(R\) 实闭，\(a,b\in R\)。判定 \(X^2+aX+b\) 在 \(R\) 上无根、处处为正的条件，并证明两者等价。
\end{exercise}
\begin{solution}
配方得 \(X^2+aX+b=(X+a/2)^2+b-a^2/4\)。若 \(a^2-4b<0\)，最后一项为正，所以多项式处处为正。若判别式为零，它在 \(-a/2\) 处为零；若判别式为正，实闭性保证它有平方根，二次公式给出两个根。因此两个所求条件都等价于 \(a^2-4b<0\)。
\end{solution}

\begin{exercise}\label{b1:ex:E-real-closure-order}
在 \(\mathbb Q(t)\) 上分别取 \(t>n\) 与 \(-t>n\) 对所有正整数 \(n\) 成立的序。证明相应的两个实闭包不可能在固定 \(\mathbb Q(t)\) 的意义下同构。
\end{exercise}
\begin{solution}
在第一个实闭包中 \(t\) 为正，故由\cref{b1:thm:real-closed-characterizations}是非零平方；在第二个中 \(t<0\)，不可能为平方。固定 \(\mathbb Q(t)\) 的域同构会保持 \(t\) 并保持平方关系，矛盾。因此实闭包唯一性中的“固定基域的序”不能略去。
\end{solution}

\begin{exercise}\label{b1:ex:E-sturm-cubic}
用 Sturm 序列求 \(f=X^3-3X+1\) 的全部实根个数，并用有理端点分别隔离它们。
\end{exercise}
\begin{solution}
取 \(f_0=X^3-3X+1\)、\(f_1=3X^2-3\)。负余数为 \(f_2=2X-1\)，再除所得负余数为 \(f_3=9/4\)。在 \(-2,-1,0,1,2\) 处，删去零项后的符号列分别为
\[
(-,+,-,+),\quad (+,-,+),\quad (+,-,-,+),\quad (-,+,+),\quad (+,+,+,+).
\]
变号数分别为 \(3,2,2,1,0\)。由\cref{b1:thm:sturm-count}，区间 \((-2,-1)\)、\((0,1)\)、\((1,2)\) 各有一个根，\((-1,0)\) 无根。已经找到三个互异根，三次多项式不再有其他根。
\end{solution}

\begin{exercise}\label{b1:ex:E-projection-square}
对实闭域 \(R\)，将集合 \(\{(x,y)\in R^2\mid y^2=x,\ y>1\}\) 投影到 \(x\) 轴，并写出不含量词的描述。
\end{exercise}
\begin{solution}
若 \(y>1\)，则 \(x=y^2>1\)。反之若 \(x>1\)，实闭性给出正平方根 \(y=\sqrt x\)。若 \(y\leq1\)，平方后会有 \(x\leq1\)，矛盾。因此投影为 \(\{x\in R\mid x>1\}\)。这一计算在所有实闭域中相同，没有使用实数上确界公理。
\end{solution}
